\documentclass[10pt,a4paper]{amsart}
\usepackage{amsmath,amssymb,amsthm,graphicx}
\usepackage{verbatim}
\newenvironment{reference}{\comment}{\endcomment}
%\newenvironment{reference}{}{}
\newenvironment{slogan}{\comment}{\endcomment}
\newenvironment{history}{\comment}{\endcomment}

% For commutative diagrams we use Xy-pic
\usepackage[all]{xy}

% We use 2cell for 2-commutative diagrams.
\xyoption{2cell}
% Theorem environments.
%
\theoremstyle{plain}
\newtheorem{theorem}[subsection]{Theorem}
\newtheorem{proposition}[subsection]{Proposition}
\newtheorem{lemma}[subsection]{Lemma}

\theoremstyle{definition}
\newtheorem{definition}[subsection]{Definition}
\newtheorem{example}[subsection]{Example}
\newtheorem{exercise}[subsection]{Exercise}
\newtheorem{situation}[subsection]{Situation}

\theoremstyle{remark}
\newtheorem{remark}[subsection]{Remark}
\newtheorem{remarks}[subsection]{Remarks}

\numberwithin{equation}{subsection}

% Macros
%
\def\lim{\mathop{\mathrm{lim}}\nolimits}
\def\colim{\mathop{\mathrm{colim}}\nolimits}
\def\Spec{\mathop{\mathrm{Spec}}}
\def\Hom{\mathop{\mathrm{Hom}}\nolimits}
\def\Ext{\mathop{\mathrm{Ext}}\nolimits}
\def\SheafHom{\mathop{\mathcal{H}\!\mathit{om}}\nolimits}
\def\SheafExt{\mathop{\mathcal{E}\!\mathit{xt}}\nolimits}
\def\Sch{\mathit{Sch}}
\def\Mor{\mathop{\mathrm{Mor}}\nolimits}
\def\Ob{\mathop{\mathrm{Ob}}\nolimits}
\def\Sh{\mathop{\mathit{Sh}}\nolimits}
\def\NL{\mathop{N\!L}\nolimits}
\def\CH{\mathop{\mathrm{CH}}\nolimits}
\def\proetale{{pro\text{-}\acute{e}tale}}
\def\etale{{\acute{e}tale}}
\def\QCoh{\mathit{QCoh}}
\def\Ker{\mathop{\mathrm{Ker}}}
\def\Im{\mathop{\mathrm{Im}}}
\def\Coker{\mathop{\mathrm{Coker}}}
\def\Coim{\mathop{\mathrm{Coim}}}

% Boxtimes
%
\DeclareMathSymbol{\boxtimes}{\mathbin}{AMSa}{"02}

%
% Macros for moduli stacks/spaces
%
\def\QCohstack{\mathcal{QC}\!\mathit{oh}}
\def\Cohstack{\mathcal{C}\!\mathit{oh}}
\def\Spacesstack{\mathcal{S}\!\mathit{paces}}
\def\Quotfunctor{\mathrm{Quot}}
\def\Hilbfunctor{\mathrm{Hilb}}
\def\Curvesstack{\mathcal{C}\!\mathit{urves}}
\def\Polarizedstack{\mathcal{P}\!\mathit{olarized}}
\def\Complexesstack{\mathcal{C}\!\mathit{omplexes}}
% \Pic is the operator that assigns to X its picard group, usage \Pic(X)
% \Picardstack_{X/B} denotes the Picard stack of X over B
% \Picardfunctor_{X/B} denotes the Picard functor of X over B
\def\Pic{\mathop{\mathrm{Pic}}\nolimits}
\def\Picardstack{\mathcal{P}\!\mathit{ic}}
\def\Picardfunctor{\mathrm{Pic}}
\def\Deformationcategory{\mathcal{D}\!\mathit{ef}}
\usepackage{cleveref}
\crefname{originalreference}{original reference}{original references}
\Crefname{originalreference}{Original reference}{Original references}
\makeatletter
\newlabel{lemma-derived-quasi-coherent-small-etale-site}{{071Q}{source}{}{sourcepointer-1}{}}
\newlabel{lemma-derived-quasi-coherent-small-etale-site@cref}{{[originalreference][0][]071Q}{[1][0][]source}{}{}{}}
\makeatother
\crefname{originallemma}{original lemma}{original lemmas}
\Crefname{originallemma}{Original lemma}{Original lemmas}
\makeatletter
\newlabel{perfect-lemma-affine-morphism-pull-push}{{08I9}{source}{}{sourcepointer-2}{}}
\newlabel{perfect-lemma-affine-morphism-pull-push@cref}{{[originallemma][0][]08I9}{[1][0][]source}{}{}{}}
\makeatother
\crefname{originalreference}{original reference}{original references}
\Crefname{originalreference}{Original reference}{Original references}
\makeatletter
\newlabel{remark-match-total-direct-images}{{08GH}{source}{}{sourcepointer-3}{}}
\newlabel{remark-match-total-direct-images@cref}{{[originalreference][0][]08GH}{[1][0][]source}{}{}{}}
\makeatother
\crefname{originallemma}{original lemma}{original lemmas}
\Crefname{originallemma}{Original lemma}{Original lemmas}
\makeatletter
\newlabel{spaces-cohomology-lemma-affine-vanishing-higher-direct-images}{{073H}{source}{}{sourcepointer-4}{}}
\newlabel{spaces-cohomology-lemma-affine-vanishing-higher-direct-images@cref}{{[originallemma][0][]073H}{[1][0][]source}{}{}{}}
\makeatother
\title{Projection formula for affine morphisms against f\_*O\_X}
\author{The Stacks Project Authors}
\date{stacks-project:spaces-perfect.tex:local-snapshot-20261012; GNU FDL 1.2 or later}
\begin{document}
\maketitle
\noindent\textit{Editorial note.} All delivered mathematical text is the author's own, copied from the pinned Stacks Project chapter file spaces-perfect.tex; no statement or proof step was rewritten, reordered or removed.\par\medskip
\noindent\textit{Editorial note.} The counted proof is the author's complete argument: the proof constructs the comparison map via the counit of the pull-push adjunction, reduces etale locally to the affine scheme case, and quotes the scheme projection formula\par\medskip
\noindent\textit{Editorial note.} The label lemma-affine-morphism-pull-push names the counted statement, and the proof environment immediately following it in the source is delivered verbatim as its proof.\par\medskip
\noindent\textit{Editorial note.} The counted statement is carried by the author's own lemma environment in the chapter source.\par\medskip
\noindent\textit{Allowed standard prerequisites (editorial).}\begin{itemize}
\item vanishing of higher direct images under affine morphisms, so f\_*O\_X -> Rf\_*O\_X is an isomorphism (Cohomology of Spaces, Lemma spaces-cohomology-lemma-affine-vanishing-higher-direct-images)
\item the canonical projection map E tensor\textasciicircum{}L Rf\_*O\_X -> Rf\_* Lf\textasciicircum{}* E adjoint to the counit Lf\textasciicircum{}*Rf\_*O\_X -> O\_X
\item the comparison of D\_QCoh on the etale site with schemes (lemma-derived-quasi-coherent-small-etale-site of the same chapter) and the compatibility of total direct images (remark-match-total-direct-images of the same chapter)
\item the scheme projection formula for affine morphisms (Derived Categories of Schemes, Lemma perfect-lemma-affine-morphism-pull-push)
\end{itemize}\medskip
\ifdefined\hypertarget\hypertarget{sourcepointer-1}{}\fi
\noindent\textit{Source reference.} Original reference 071Q: sources/stacks-spaces-perfect/pinned-original-source.tex, source lines 402--402. Prerequisite pointer invoked at proof line(s) 1011; it names a standard fact of the same chapter that the counted argument uses, and every step of the counted proof built on it is retained verbatim in the delivered body.\par\medskip
\ifdefined\hypertarget\hypertarget{sourcepointer-2}{}\fi
\noindent\textit{Source reference.} Original lemma 08I9: sources/stacks-chow/tags.txt, source lines 10807--10807. Prerequisite pointer invoked at proof line(s) 1010; it names a standard fact of a sibling chapter that the counted argument uses, and every step of the counted proof built on it is retained verbatim in the delivered body.\par\medskip
\ifdefined\hypertarget\hypertarget{sourcepointer-3}{}\fi
\noindent\textit{Source reference.} Original reference 08GH: sources/stacks-spaces-perfect/pinned-original-source.tex, source lines 896--896. Prerequisite pointer invoked at proof line(s) 1012; it names a standard fact of the same chapter that the counted argument uses, and every step of the counted proof built on it is retained verbatim in the delivered body.\par\medskip
\ifdefined\hypertarget\hypertarget{sourcepointer-4}{}\fi
\noindent\textit{Source reference.} Original lemma 073H: sources/stacks-chow/tags.txt, source lines 9061--9061. Prerequisite pointer invoked at proof line(s) 994; it names a standard fact of a sibling chapter that the counted argument uses, and every step of the counted proof built on it is retained verbatim in the delivered body.\par\medskip
\begin{lemma}
\label{lemma-affine-morphism-pull-push}
Let $S$ be a scheme. Let $f : X \to Y$ be an affine morphism of algebraic
spaces over $S$. For $E$ in $D_\QCoh(\mathcal{O}_Y)$ we have
$Rf_* Lf^* E = E \otimes^\mathbf{L}_{\mathcal{O}_Y} f_*\mathcal{O}_X$.
\end{lemma}

\begin{proof}
Since $f$ is affine the map $f_*\mathcal{O}_X \to Rf_*\mathcal{O}_X$
is an isomorphism (Cohomology of Spaces, Lemma
\ref{spaces-cohomology-lemma-affine-vanishing-higher-direct-images}).
There is a canonical map
$E \otimes^\mathbf{L} f_*\mathcal{O}_X =
E \otimes^\mathbf{L} Rf_*\mathcal{O}_X \to Rf_* Lf^* E$
adjoint to the map
$$
Lf^*(E \otimes^\mathbf{L} Rf_*\mathcal{O}_X) =
Lf^*E \otimes^\mathbf{L} Lf^*Rf_*\mathcal{O}_X \longrightarrow
Lf^* E \otimes^\mathbf{L} \mathcal{O}_X = Lf^* E
$$
coming from $1 : Lf^*E \to Lf^*E$ and the canonical map
$Lf^*Rf_*\mathcal{O}_X \to \mathcal{O}_X$. To check the map so constructed
is an isomorphism we may work locally on $Y$. Hence we may assume
$Y$ and therefore $X$ is affine. In this case the problem reduces to the
case of schemes
(Derived Categories of Schemes, Lemma
\ref{perfect-lemma-affine-morphism-pull-push})
via Lemma \ref{lemma-derived-quasi-coherent-small-etale-site} and
Remark \ref{remark-match-total-direct-images}.
\end{proof}

\end{document}
